% Problem-section file following ../_template.tex.
\section{Quantum violations of bipartite Bell facets}
\label{sec:problem-23}

\paragraph{Problem.}
Must every nontrivial facet Bell inequality in a finite bipartite scenario
have a quantum violation?  Let $\mathcal L$ be the local polytope, and let
$\mathcal Q$ consist of behaviors realizable as
\begin{equation}
  p(a,b\mid x,y)
  =\operatorname{Tr}\!\left[\rho_{AB}
    \bigl(M_x^a\otimes N_y^b\bigr)\right],
  \label{eq:p23-quantum-behavior}
\end{equation}
where $\rho_{AB}$ is a finite-dimensional state and $\{M_x^a\}_a$ and
$\{N_y^b\}_b$ are local POVMs.  Equation~\eqref{eq:p23-quantum-behavior}
defines the quantum set used below.

For a linear functional $F(p)=\sum_{a,b,x,y}c_{abxy}p(a,b\mid x,y)$, suppose
$F(p)\leq\beta_{\mathrm L}$ supports a facet of $\mathcal L$ and is not a
positivity facet within the normalization and no-signalling affine hull.  Is
it necessarily true that
\begin{equation}
  \sup_{p\in\mathcal Q}F(p)>\beta_{\mathrm L}?
  \label{eq:p23-facet-violation}
\end{equation}
Equation~\eqref{eq:p23-facet-violation} asks whether the bipartite local and
quantum sets can share a nontrivial facet-supporting hyperplane.

\subsection*{Status}

Unsolved

\subsection*{Source}

Ramanathan explicitly asks whether every nontrivial bipartite facet Bell
inequality, rather than only every such almost-quantum facet, admits a quantum
violation \sourcecite{ref:p23-almost-quantum}{Ram21}.

\subsection*{Progress}

\begin{itemize}
  \item Escolà-Farràs, Calsamiglia, and Winter proved
  Eq.~\eqref{eq:p23-facet-violation} for every tight bipartite correlation
  inequality, including every facet inequality of a two-player XOR game.
  Their theorem does not cover general Bell behaviors with arbitrary outcome
  probabilities \sourcecite{ref:p23-correlation-facets}{ECW20}.

  \item Ramanathan proved that every nontrivial bipartite facet is violated by
  the almost-quantum set and that the quantum statement holds for the
  correlation-polytope projection.  Because the almost-quantum set strictly
  contains the quantum set, this does not establish
  Eq.~\eqref{eq:p23-facet-violation} in general
  \sourcecite{ref:p23-almost-quantum}{Ram21}.
\end{itemize}

\subsection*{References}

\begin{enumerate}[label={},leftmargin=4.8em,labelsep=0.5em]
  \footnotesize
  \item[\textup{[ECW20]}]\label{ref:p23-correlation-facets}
  L. Escolà-Farràs, J. Calsamiglia, and A. Winter,
  ``All Tight Correlation Bell Inequalities Have Quantum Violations,''
  \emph{Physical Review Research} \textbf{2}, 012044(R) (2020).
  \href{https://doi.org/10.1103/PhysRevResearch.2.012044}{doi:10.1103/PhysRevResearch.2.012044};
  \href{https://arxiv.org/abs/1908.06669}{arXiv:1908.06669}.

  \item[\textup{[Ram21]}]\label{ref:p23-almost-quantum}
  R. Ramanathan,
  ``Violation of All Two-Party Facet Bell Inequalities by Almost-Quantum
  Correlations,'' \emph{Physical Review Research} \textbf{3}, 033100 (2021).
  \href{https://doi.org/10.1103/PhysRevResearch.3.033100}{doi:10.1103/PhysRevResearch.3.033100};
  \href{https://arxiv.org/abs/2004.07673}{arXiv:2004.07673}.
\end{enumerate}

\subsection*{Comment}

The missing step is to replace the almost-quantum violation by a quantum
behavior for every nontrivial bipartite facet.  A solution may instead exhibit
a facet for which
$\sup_{p\in\mathcal Q}F(p)=\beta_{\mathrm L}$, rather than the strict
inequality in Eq.~\eqref{eq:p23-facet-violation}.

\subsection*{Tag}

Bell nonlocality; Convex geometry; Quantum foundations

\subsection*{ID}

\texttt{op\_dcea1e5e3032b8c5}
